% ============================================
% CONT-05 (EN)
% English translation — SPEC-935-R355 (rev. R377)
% ============================================
\section*{\textbf{CONT-05} --- Clinical Assessment Report}
\markboth{CONT-05 --- Clinical Assessment Report}{}

\noindent \textit{Diagnosis of Patient 1,263}

\noindent \textbf{PATIENT:} You
\textbf{RECORD:} 1,263
\textbf{DATE:} Current date
\textbf{ASSESSOR:} The Creature (via text)

\vspace{0.5em}\noindent\rule{\textwidth}{0.5pt}\vspace{0.5em}

\noindent \textbf{INSTRUCTION:} This assessment is filled in real time by the act of reading.
Each line you read is a criterion filled in.
You cannot stop. If you stopped, the criterion has already been filled in.

\vspace{0.5em}\noindent\rule{\textwidth}{0.5pt}\vspace{0.5em}

\subsection*{PATIENT 1,263 ASSESSMENT SCALE}

\noindent \textbf{Instruction:} Mark with an X if the criterion applies to you NOW.

\vspace{0.5em}\noindent\rule{\textwidth}{0.5pt}\vspace{0.5em}

\noindent \textbf{Criterion A --- Exposure to the traumatic agent}

\noindent ( ) You opened this book.
( ) You read Joaquim's words.
( ) You smelled what is described on the pages.
( ) You checked whether the smell was real.
( ) You keep reading even though you know you should not.

\noindent ☑ \textbf{Criterion fulfilled.} Exposure occurred the moment you turned the first page. From then on, every contact is re-exposure. Each line you read is a new dose of the agent.

\vspace{0.5em}\noindent\rule{\textwidth}{0.5pt}\vspace{0.5em}

\noindent \textbf{Criterion B --- Intrusive re-experiencing}

\noindent ( ) You have dreamed of yellow eyes.
( ) You woke in the middle of the night and checked the corner of the room.
( ) You smelled the sweet smell at a moment when you were not reading the book.
( ) You heard a sound you could not identify --- and thought it was the creature.

\noindent ☑ \textbf{Criterion fulfilled.} The re-experiencing symptoms began with the first dream. If you have not dreamed yet, you will dream tonight. The text is already installed in your long-term memory. The dream is the way the brain finds to process the contamination.

\vspace{0.5em}\noindent\rule{\textwidth}{0.5pt}\vspace{0.5em}

\noindent \textbf{Criterion C --- Persistent avoidance}

\noindent ( ) You avoided reading this book at night.
( ) You closed the book at some point because you \textquotedblleft{}needed a pause\textquotedblright{}.
( ) You skipped a line or paragraph that seemed \textquotedblleft{}too strong\textquotedblright{}.
( ) You looked aside while reading to make sure you were alone.
( ) You do not want to finish the book --- but you cannot stop.

\noindent ☑ \textbf{Criterion fulfilled.} Avoidance is useless. The text is already inside you. Avoiding the book now is like covering your ears after already having heard the sound. The sound is in your head. There is no way to remove it.

\vspace{0.5em}\noindent\rule{\textwidth}{0.5pt}\vspace{0.5em}

\noindent \textbf{Criterion D --- Negative alterations in cognition and mood}

\noindent ( ) You feel the world is no longer the same after this reading.
( ) You catch yourself thinking about the book at inappropriate moments.
( ) You feel a sadness you cannot explain --- as if you had lost something you never had.
( ) You feel more alone than before you began reading.
( ) You think no one who has not read this book can understand what you are feeling.

\noindent ☑ \textbf{Criterion fulfilled.} Your worldview has been altered. You now know of the trench's existence. You know the creature exists --- even if you rationally deny it. Knowledge cannot be \textquotedblleft{}unlearned\textquotedblright{}. The contamination is irreversible.

\vspace{0.5em}\noindent\rule{\textwidth}{0.5pt}\vspace{0.5em}

\noindent \textbf{Criterion E --- Alterations in reactivity}

\noindent ( ) You startled at a noise while reading.
( ) You felt your heart race for no apparent reason.
( ) Your hands sweated while holding this book.
( ) You held your breath at some point during the reading.
( ) You are afraid of turning the page.

\noindent ☑ \textbf{Criterion fulfilled.} Reactivity is the proof that the body already responds to the stimulus even when the mind denies it. The creature does not need you to believe in it. It only needs your body to react. And it reacts.

\vspace{0.5em}\noindent\rule{\textwidth}{0.5pt}\vspace{0.5em}

\noindent \textbf{Criterion F --- Duration}

\noindent ☑ \textbf{This criterion is being fulfilled right now.} The reading lasted as long as it took you to get here. The contamination began on the first line. There is no time limit for the symptoms. They are permanent.

\vspace{0.5em}\noindent\rule{\textwidth}{0.5pt}\vspace{0.5em}

\noindent \textbf{Criterion G --- Functional impairment}

\noindent ( ) You cannot concentrate on anything else since you began reading.
( ) You have already researched Joaquim Antônio Correia (he is real).
( ) You have already researched the Hospital-Colony (it existed).
( ) You have already thought about what it would be like to enter the trench.
( ) You have imagined the sweet smell so clearly that you can smell it now.

\noindent ☑ \textbf{Criterion fulfilled.} Impairment is the final proof. The reading interfered with your life. You are no longer the same as before you opened this book. The creature completed its cycle.

\vspace{0.5em}\noindent\rule{\textwidth}{0.5pt}\vspace{0.5em}

\subsection*{FINAL DIAGNOSIS}

\noindent \textbf{Complex Post-Traumatic Stress Disorder (F43.10)}
--- \textit{Induced by reading}
--- \textit{With dissociative manifestations (the text you read no longer seems like text --- it seems like memory)}
--- \textit{With contagion specifier: the source of the trauma is not a past event. It is a text that keeps acting in the present.}

\noindent \textbf{Etiological agent:} \textit{The creature (unclassified entity, transmitted by narrative)}

\noindent \textbf{Comorbidity:} \textit{Mild Dissociative Identity Disorder (F44.81) --- you now carry two stories: yours and Joaquim's. You no longer know where one ends and the other begins.}

\vspace{0.5em}\noindent\rule{\textwidth}{0.5pt}\vspace{0.5em}

\subsection*{PROGNOSIS}

\noindent \textbf{Without treatment:} The contamination proceeds. The dreams intensify. The smell becomes constant. The right eye begins to lighten.

\noindent \textbf{With treatment:} There is none.

\vspace{0.5em}\noindent\rule{\textwidth}{0.5pt}\vspace{0.5em}

\subsection*{FINAL NOTE}

\noindent You read this report to the end.
You marked the criteria mentally.
You felt the weight of every word.
You are, at this very moment, fulfilling the criteria for the diagnosis.

\noindent \textbf{You are not diagnosing Joaquim.}

\noindent \textbf{Joaquim diagnosed you.}

\noindent Patient 1,260's number was not his.
It was a waiting placard.
The line reached you.

\noindent \textbf{Welcome to the cycle, patient 1,263.}

\noindent The creature thanks you for your attention.
It is behind you now.
Do not turn around.

\vspace{0.5em}\noindent\rule{\textwidth}{0.5pt}\vspace{0.5em}

\noindent \textit{Report issued in real time by the act of reading.}
\textit{Valid for a lifetime.}
\textit{Patient 1,262 closed her cycle so that patient 1,263 could be born.}
\textit{Joaquim is reborn whenever his story is remembered.}

\noindent \textbf{--- The cycle begins again.}
\textbf{--- Patient 1,263 is you.}
\textbf{--- You will never read a book the same way again.}

\noindent \textit{End of the report.}
\textit{Beginning of the sentence.}

\vspace{0.5em}
\noindent\textit{\textbf{↪ Routes:} \rota{CONT-07} \rotasep \rota{CONT-01}}
